\documentclass{acm_proc_article-sp}
\usepackage[utf8]{inputenc}

\def\inputfig#1{\input #1}
\def\inputtex#1{\input #1}
\def\inputal#1{\input #1}
\def\inputcode#1{\input #1}

\inputtex{logos.tex}
\inputtex{refmacros.tex}
\inputtex{other-macros.tex}

\begin{document}
\title{Highly configurable Common Lisp reader}
\numberofauthors{2}
\author{\alignauthor
Jan Morningen\\
\affaddr{xxx}\\
\affaddr{xxx}\\
\affaddr{xxx}\\
\email{scymtym@gmx.net}}
\author{\alignauthor
Robert Strandh\\
\affaddr{Unaffiliated}\\
\email{robert.strandh@gmail.com}}

\toappear{Permission to make digital or hard copies of all or part of
  this work for personal or classroom use is granted without fee
  provided that copies are not made or distributed for profit or
  commercial advantage and that copies bear this notice and the full
  citation on the first page. Copyrights for components of this work
  owned by others than the author(s) must be honored. Abstracting with
  credit is permitted. To copy otherwise, or republish, to post on
  servers or to redistribute to lists, requires prior specific
  permission and/or a fee. Request permissions from
  Permissions@acm.org.

  ELS '26, August 14 - 17 2014, Montreal, QC, Canada
  Copyright is held by the owner/author(s). Publication rights licensed to ACM.
  ACM 978-1-4503-2931-6/14/08\$15.00.
  http://dx.doi.org/10.1145/2635648.2635656}

\maketitle

\begin{abstract}
The \cl{} \texttt{read} function (usually called the \cl{}
\emph{reader}) is an essential part of every \cl{} implementation,
because it used to read \cl{} code for subsequent evaluation.  The
readers of most existing \cl{} implementations do this well, but
little more.  However, there are many situations where it is useful to
read \cl{} code for other purposes as well, and in those situations it
can be essential to configure the reader in various ways, for instance
so that it does not fail for symbols with incorrect package markers.
In this paper, we describe \eclector{}, a highly configurable \cl{}
reader.  Configuration is accomplished by the extensive use of \cl{}
generic functions called by \eclector{} and that can be specialized by
client code.
\end{abstract}

\category{D.3.4}{Programming Languages}{Processors}
[Code generation, Run-time environments]

\terms{Algorithms, Languages}

\keywords{\clos{}, \cl{}, Parsing}

\inputtex{sec-introduction.tex}
\inputtex{sec-previous.tex}
\inputtex{sec-our-method.tex}
\inputtex{sec-conclusions.tex}

\bibliographystyle{abbrv}
\bibliography{eclector}
\end{document}
